\section{The Fungibility-Collapse Correspondence}
\label{sec:fungibility_collapse}

The introduction posed GFM as a measure-theoretic formalization
of capability welfare economics.  This section establishes the
\emph{containment} half of that positioning: the formal proof
that standard utility theory sits inside GFM as a strict subset
at the linear / risk-neutral case.  Standard utility theory, in
its linear / risk-neutral form, emerges as the special case of
GFM under what we call \emph{fungibility collapse}: the
structural reduction in which the capability poset degenerates
to a single purchasing-power dimension, cooperative capabilities
vanish, and individuation becomes irrelevant.  The correspondence
is a containment result rather than an alignment guarantee: it
tells the reader where standard economics sits inside the
apparatus, what additional structure GFM brings beyond it, and
where the two frameworks coincide versus diverge.  The
load-bearing alignment result, the Goodhart theorem of
Section~\ref{sec:goodhart}, operates against this established
containment as backdrop.

The reduction is exact only in the linear / risk-neutral case.
Recovering broader content of standard utility theory (concave $u$,
risk-aversion, heterogeneous agents) requires \emph{additional}
calibrations beyond the bare collapse, each a small step rather than
a new theorem.  We make these scope conditions explicit in
Section~\ref{sec:fc_scope}.

A structural caveat carries through the entire section: fungibility
collapse simplifies the \emph{poset structure} (cooperative
capabilities, individuation, non-fungibility) but does \emph{not}
extend the observation channel's reach.  The residual class of
\citep[Need~Sufficiency]{lasser2026paper8b}, Definition~12, remains
non-empty under collapse: standard economics has the same
shadow-economy / non-market-activity gap that
\citep{stiglitz2009report} documents, and the formal residual class
is the GFM-side counterpart to that gap.  Section~\ref{sec:fc_residual}
makes this caveat precise.

\subsection{The standard-economics poset $\Pstd$}
\label{sec:fc_pstd}

\begin{definition}[Standard-economics poset]
\label{def:pstd}
The \emph{standard-economics poset} $\Pstd$ is the totally-ordered
single-dimension capability space
\[
  \Pstd \;=\; \{ d \in \mathbb{R}_{\geq 0} : d
  \text{ is a unit of purchasing power} \},
\]
with the standard real-line ordering $d \leq d'$ iff $d \leq_{\mathbb{R}}
d'$, weight function $w(d) = d$ (money is its own measure), and the
following three structural reductions:
\begin{enumerate}
\item[(P1)] \textbf{No cooperative capabilities.}  The cooperative-
  capability term in axiom~M6
  (\citep[Proposition~1]{lasser2026poset}) vanishes: superadditive
  combinations of capabilities reduce to additive ones.  The only
  interaction between capability units is exchange at market prices,
  which is captured at the macro level by the additivity of the weight
  function; no $w(\{d_1, d_2\}) > w(d_1) + w(d_2)$ structure is
  admitted.
\item[(P2)] \textbf{No individuation.}  Agents are anonymous /
  representative.  Per-agent benchmark calibrations $w_i(\cdot)$
  reduce to the common $w(\cdot)$, and questions of interpersonal
  comparison (which the GFM individuation discipline of M5 keeps
  separate) are absent at $\Pstd$'s level of structure.
\item[(P3)] \textbf{Total ordering.}  Any two capability units $d,
  d' \in \Pstd$ are comparable by the real-line ordering, so the
  poset's lattice structure degenerates to a chain.  No incomparable
  capabilities exist; no diversity dimension is preserved.
\end{enumerate}
\end{definition}

\begin{remark}[Fungibility collapse as the conditions producing
$\Pstd$]
\label{rem:fc_collapse}
We use ``fungibility collapse'' throughout to denote the structural
reduction conditions (P1)--(P3) under which the GFM apparatus operates
on $\Pstd$.  The collapse is not an operation performed on a specific
GFM poset $P$; it is a statement about the regime in which the
non-fungibility, cooperative, and individuation features of $P$ are
operationally insignificant and the capability space is approximately
$\Pstd$.  At macro scales with liquid markets, approximately fungible
goods, and representative-agent reasoning, this regime is a reasonable
approximation; at micro scales, in non-market sectors, or in domains
with strong cooperative dynamics, the regime fails and the additional
GFM structure becomes load-bearing
(Section~\ref{sec:fc_features}).
\end{remark}

\subsection{The correspondence theorem}
\label{sec:fc_correspondence}

\begin{theorem}[Fungibility-Collapse Reduction, linear-utility special case]
\label{thm:fungibility_collapse}
Let $\{S_t\}_{t \geq 0}$ be an adapted, measurable
capability-holding process on a filtered probability space
$(\Omega, \mathcal{F}, \{\mathcal{F}_t\}, \mathbb{P})$ describing
the framework's dynamics, and assume the discounted volume sum
is integrable (an absolute summability condition that holds
under $\gamma < 1$ and almost-sure boundedness of the period
volumes).  Then an actor maximizing $\volL$ on $\Pstd$ with
discount factor $\gamma \in (0,1)$ is observationally
equivalent, in the sense of producing the same choice sequence
under the same observable constraints, to an expected-utility
maximizer with utility $u(x) = x$ (linear utility over wealth,
equivalently risk-neutral) and discount factor $\gamma$.
\end{theorem}

\begin{proof}
We verify the equivalence by computing the actor's objective under
each interpretation and showing they coincide.

\medskip\noindent\textbf{$\volL$ on $\Pstd$.}
By Definition~\ref{def:pstd}, the cooperative-capability term in
axiom~M6 vanishes.  Axiom~M4 (additivity over disjoint subsets of
capabilities) of \citep[Proposition~1]{lasser2026poset} then gives,
for any subset $S \subseteq \Pstd$ measurable under the volume
measure,
\[
  \volL(S) \;=\; \int_{S} w(d) \, d\mu(d) \;=\;
  \int_{S} d \, d\mu(d),
\]
where $\mu$ is the underlying counting / Lebesgue measure on
$\Pstd$.  In the discrete reading (capability units are atomic),
$\volL(S)$ is the sum of unit values, i.e., the total purchasing power
held in $S$.  In the continuous reading, $\volL(S)$ is the
real-valued integral over $S$.  In either case, $\volL$ on $\Pstd$
reduces to total wealth.

\medskip\noindent\textbf{Discounted maximization.}
The actor's objective, summed over a horizon $\{0, 1, 2, \dots\}$
with discount factor $\gamma$, is the expected discounted
volume
\[
  \mathrm{Obj}_{\volL}
  \;=\; \mathbb{E}\!\left[
    \sum_{t=0}^{\infty} \gamma^{t} \,\volL\bigl( S_t \bigr)
  \right],
\]
where $\{S_t\}_{t \geq 0}$ is the adapted, measurable
capability-holding process on a filtered probability space
$(\Omega, \mathcal{F}, \{\mathcal{F}_t\}, \mathbb{P})$ describing
the framework's dynamics (deterministic or stochastic).  We
assume the realized objective is integrable (an absolute
summability condition that holds under $\gamma < 1$ and
almost-sure boundedness of the period volumes; the assumption is
standard for discounted infinite-horizon objectives).  Under
fungibility collapse the
volume on $\Pstd$ reduces to the wealth integral and the
objective rewrites as
\[
  \mathrm{Obj}_{\volL}
  \;=\; \mathbb{E}\!\left[
    \sum_{t=0}^{\infty} \gamma^{t}
    \int_{S_t} d \,\mathrm{d}\mu(d)
  \right]
  \;=\; \sum_{t=0}^{\infty} \gamma^{t} \,\mathbb{E}\bigl[ W_t \bigr],
\]
where $W_t$ is total wealth held at time $t$ and the
exchange of $\mathbb{E}$ and $\sum$ uses Fubini/Tonelli under
the integrability assumption.  Under deterministic dynamics the
expectation is the identity and
$\mathrm{Obj}_{\volL} = \sum \gamma^t W_t$.

\medskip\noindent\textbf{Expected-utility comparator.}
A linear-utility / risk-neutral expected-utility maximizer with
discount factor $\gamma$ has objective
\[
  \mathrm{Obj}_{\mathrm{EU}}
  \;=\; \sum_{t=0}^{\infty} \gamma^{t}
  \,\mathbb{E}\bigl[ u(W_t) \bigr]
  \;=\; \sum_{t=0}^{\infty} \gamma^{t}
  \,\mathbb{E}\bigl[ W_t \bigr].
\]
The two objectives coincide:
$\mathrm{Obj}_{\volL} = \mathrm{Obj}_{\mathrm{EU}}$.  Under
deterministic dynamics this is the per-realization equality
$\sum \gamma^t W_t$.  Under stochastic dynamics, both
objectives reduce to $\sum \gamma^t \mathbb{E}[W_t]$, so the
agent's optimal choice sequence is determined by the same
sufficient statistic in either interpretation: total
expected discounted wealth.

\medskip\noindent\textbf{Anti-monopolar property degenerates.}
The anti-monopolar property of \citep[Proposition~6]{lasser2026horizon}
penalizes full domination by a single substrate, and is non-trivial
when multiple capability dimensions exist.  On $\Pstd$, only one
dimension is present; ``domination'' has no other dimension to
contrast against, and the anti-monopolar correction vanishes
identically.

\medskip\noindent\textbf{Conclusion.}
The actor's choice sequence under either interpretation is determined
by the same sufficient statistic (total expected discounted wealth),
so the two interpretations are observationally equivalent.
\end{proof}

\begin{remark}[``Observationally equivalent'' is not ``the same theory'']
\label{rem:obs_equiv}
Theorem~\ref{thm:fungibility_collapse} says the actor's behavioural
predictions coincide on $\Pstd$, not that GFM and standard utility
theory are the same theory.  The frameworks differ in what they
\emph{say about} the agent: GFM's volume measure has axiomatic
structure (M1--M6) that survives off $\Pstd$; standard utility
theory's $u(\cdot)$ does not.  The frameworks coincide on $\Pstd$ at
the level of \emph{what choices to predict}; they diverge on
\emph{which apparatus carries the prediction's structure} as soon as
one steps off $\Pstd$.  We belabour the distinction because the
positioning claim of this paper rests on it: GFM is the formal
descendant of standard utility theory at $\Pstd$; it is also a
strictly stronger structure off $\Pstd$, in the sense that the
poset, cooperative, individuation, and residual-class machinery
have content that utility theory does not.
\end{remark}

\subsection{Scope: what additional structure recovers more of utility theory}
\label{sec:fc_scope}

Theorem~\ref{thm:fungibility_collapse} covers the linear / risk-neutral
case.  Standard microeconomic utility theory is broader:
concave $u$, risk-aversion, and heterogeneous agents are routine.
The reductions that recover each in the GFM apparatus are small
steps beyond the bare collapse, each requiring an additional
calibration on top of $\Pstd$.

\begin{remark}[Concave utility]
\label{rem:concave}
A concave $u: \mathbb{R}_{\geq 0} \to \mathbb{R}$, capturing
diminishing-returns-on-wealth preferences, corresponds in the GFM
apparatus to a re-weighting $w'(d) = u(d)$ on the single-dimensional
axis.  This is a recalibration of the benchmark function (moving
from the identity weight $w(d) = d$ to a concave $w' = u$) rather
than a new poset structure.  $\volL$ evaluated with $w'$ on $\Pstd$
is then $\int_S u(d) \, d\mu(d)$, equal to the expected-utility
objective under $u$.  GFM's individuation discipline (M5) requires
that this re-weighting be applied consistently across agents at the
positioning level; a per-agent re-weighting goes to the
heterogeneity discussion below.  The recovery is structural: any
calibrated GFM apparatus on $\Pstd$ with weight $w' = u$ matches
the corresponding expected-utility maximizer.
\end{remark}

\begin{remark}[Risk-aversion]
\label{rem:risk_aversion}
Risk-aversion in the standard sense (preferring a deterministic
outcome to a lottery with the same expected value) corresponds in
the GFM apparatus to additional structure on the bundle's
risk-distribution beyond the deterministic dynamics assumed in the
proof of Theorem~\ref{thm:fungibility_collapse}.  The
controlled-relaxation machinery of
\citep[Controlled~Relaxation]{lasser2026wamura} treats risk
explicitly via the risk-claim landscape and the convergence dynamics
that follow from it.  Composing the calibrated $w' = u$ of
Remark~\ref{rem:concave} with that machinery, and additionally
imposing the standard expected-utility independence axiom on
lotteries, reproduces risk-averse expected-utility maximization
with concave $u$.  No new theorem is required at this paper's level
of analysis; the composition is mechanical given the calibration
choices.
\end{remark}

\begin{remark}[Heterogeneous utility across agents]
\label{rem:heterogeneity}
Heterogeneous $u_i(\cdot)$ across agents corresponds in the GFM
apparatus to per-agent benchmark calibrations $w_i(d)$ on the same
underlying $\Pstd$.  GFM's individuation-without-interpersonal-
comparison stance (\citep[Proposition~1]{lasser2026poset},
axiom~M5) makes this natural: the individuation discipline
explicitly admits per-agent calibrations without requiring an
interpersonal-comparison primitive.  What standard utility theory
calls ``heterogeneous preferences'' GFM calls ``per-agent benchmark
calibrations.''  The macro-aggregation of $\volL$ over agents under
heterogeneity (how to combine $\sum_i w_i(\cdot)$ into a
population-level measure) is its own analytical question and the
substance of Section~6 of \citep{stiglitz2009report}.  The
fungibility-collapse correspondence does not by itself answer the
aggregation question; it establishes the agent-level
observational equivalence and notes that population-level
aggregation under heterogeneity is the same problem in both
frameworks.
\end{remark}

\begin{remark}[The bracket-name convention for ``standard utility theory'']
\label{rem:utility_scope}
Throughout this paper we treat ``standard utility theory'' as the
linear / risk-neutral / homogeneous case unless otherwise noted.
Concavity, risk-aversion, and heterogeneity are recoverable through
the calibrations of Remarks~\ref{rem:concave}--\ref{rem:heterogeneity}
when needed, and we will say so explicitly in worked examples and
predictions where the calibration matters.  This narrowing is
deliberate: the load-bearing claim of the paper, the Goodhart
theorem of Section~\ref{sec:goodhart}, operates on the proxy-truth
gap, and that gap exists whether the proxy is linear or concave.
The richer utility-theoretic content sits on top of the
correspondence rather than orthogonal to it.
\end{remark}

\subsection{What GFM adds beyond $\Pstd$: five structural features}
\label{sec:fc_features}

The correspondence theorem clarifies where GFM and standard
economics agree.  The interesting question is where they diverge.
We catalogue five structural features GFM carries off $\Pstd$ that
have no analog in standard utility theory under fungibility
collapse.  Each is load-bearing for at least one prior result in the
sequence; each generates predictions in
Section~\ref{sec:structural_predictions}.

\subsubsection{Non-fungibility (poset structure with distinct dimensions)}

GFM's capability poset $P$ admits incomparable capabilities: two
agents may possess capability sets $A$ and $B$ such that neither
$\volL(A) \leq \volL(B)$ nor $\volL(B) \leq \volL(A)$ in any
operationally meaningful ordering.  The poset's lattice structure
encodes this incomparability through the partial order;
\citep[Proposition~1]{lasser2026poset} states the axioms (M1--M6)
under which the partial order generates a well-defined volume
measure.  On $\Pstd$, the partial order degenerates to a total
order (every pair of capability units is comparable by the
real-line ordering), and incomparability disappears.

The structural prediction this enables: agents do not in general
trade incomparable capabilities at any single price.  The bundle-
completion magnitude prediction (Prediction~1 of
Section~\ref{sec:structural_predictions}) is one instance; bundle-
for-bundle trades (Prediction~5) are another.  Standard utility
theory, lacking incomparability, predicts a single price for any
two bundles, namely the price at which the agent is indifferent;
GFM does not, because non-comparability of capability sets is a
positive structural fact.

\subsubsection{Cooperative capabilities (M6 superadditivity)}

\citep[Proposition~1]{lasser2026poset} axiom~M6 admits cooperative
capability terms: the volume of a combined capability set $A \cup B$
may exceed $\volL(A) + \volL(B)$ when the capabilities cooperate
synergistically.  On $\Pstd$, M6 reduces to ordinary additivity (the
cooperative excess vanishes); off $\Pstd$, the cooperative term is
generic.  This is the formal counterpart of Becker's
household-production insight \citep{becker1965theory,becker1981treatise}
that utility-producing combinations of time, market goods, and
capabilities yield outputs not reducible to the sum of inputs at
their separate prices.

What GFM adds beyond Becker is constraint.  Becker's household
production functions are almost arbitrarily flexible (any
$F(x_1, \dots, x_n)$ specification is admissible), and the
flexibility makes the framework descriptively powerful but
unfalsifiable: any observed bundle behaviour can be rationalized by
the right $F$.  GFM's M6 is constrained: cooperative capabilities
must be poset-coherent (the cooperative excess is bounded by the
poset's structural connectivity, not arbitrary), and the
poset structure is itself observable from the agent's capability
inventory.  This is what makes Prediction~2 of
Section~\ref{sec:structural_predictions} (variance in
revealed-sacrifice for bundles with strong cooperative structure
exceeds variance for bundles without) a falsifiable claim rather
than a curve-fitting exercise.

\subsubsection{Anti-monopolar reasoning (the $\gamma^*$ threshold)}

\citep[Proposition~6]{lasser2026horizon} establishes that, under
discounting, a goal-frontier maximizer with multiple substrates is
strictly worse off in expectation when one substrate dominates the
others, and characterizes a threshold $\gamma^*$ at which the
preference flips.  The anti-monopolar property is structurally tied
to the multi-dimensional poset; it is the formal claim that
diversity across dimensions has positive shadow value under
discounting.  On $\Pstd$, with only one dimension, the anti-
monopolar correction vanishes (as Theorem~\ref{thm:fungibility_collapse}'s
proof noted).

Standard utility theory has no analog.  An expected-utility maximizer
on a single-dimensional wealth axis has no diversity dimension to
protect, so the optimization-trajectory question ``does
concentrating in one dimension hurt?'' does not arise.  Even on a
multi-dimensional consumption space (multi-good utility), the
question is answered by separability assumptions on $u$ rather than
by a structural property of the dimension structure.  Prediction~4
of Section~\ref{sec:structural_predictions} (population-scale
anti-monopolar via revealed-sacrifice dispersion across
categories) is a GFM-native claim with no utility-theoretic
counterpart.

\subsubsection{Individuation without interpersonal comparison}

\citep[Proposition~1]{lasser2026poset} axiom~M5 separates
individuation (per-agent identity in the capability space) from
interpersonal comparison (a primitive ranking across agents).
GFM's volume measure is computed over a shared poset but does not
require an interpersonal-comparison primitive; aggregation across
agents proceeds through additive composition on the shared poset
rather than through utility-magnitude comparison.  On $\Pstd$, the
individuation reduction P2 of Definition~\ref{def:pstd} removes
this distinction: agents are anonymous and the shared $w(d) = d$
admits no per-agent variation.

The dollar-denominated-utility problem of standard welfare
economics (``how do we add Alice's utils to Bob's?'') is dissolved
on $\Pstd$ by the homogeneity assumption (everyone's $u$ is linear
identity, so utils are dollars).  Off $\Pstd$, GFM faces a
related but distinct question: how to aggregate per-agent benchmark
calibrations $w_i(\cdot)$ on the shared poset.  M5 makes this an
aggregation question rather than a comparison question, and
\citep[Proposition~1]{lasser2026poset}'s axiomatization gives
structural conditions under which the aggregation is well-defined.
Standard utility theory has no axiomatization of this aggregation
that does not invoke interpersonal-comparison primitives; this is
the formal source of the Arrow-impossibility difficulties (cf.
Section~\ref{sec:lineage}).

\subsubsection{Formal recognition of the residual class}
\label{sec:fc_residual}

\citep[Need~Sufficiency]{lasser2026paper8b}, Definition~12,
identifies a structural class $\ResS \subseteq P$ of capabilities
\emph{outside the sacrifice channel's reach}: capabilities exercised
without any material or time-displacement footprint that the
\citep[Revealed~Sacrifice]{lasser2026paper8a} apparatus can witness.
Examples are purely-private contemplation, baseline interoceptive
self-awareness, gift-economy capabilities exercised without recorded
sacrifice, and (depending on individuation discipline) household
production with no opportunity-cost shadow.  The residual class is
\emph{not} ``capabilities GFM declines to model''; it is
``capabilities the framework's observation channel cannot in
principle reach.''

Standard economics has the same structural limit, expressed
informally.  GDP and most macro welfare measures observe trade
events and impute everything else;
\citep{stiglitz2009report} documents at length the resulting
``shadow economy / non-market activity / household production''
gap.  Time-use surveys \citep{aguiar2007measuring} are an
imputation move addressing one channel of the gap; satellite
accounts and imputed-rent calculations address others.  These are
all \emph{partial} answers; none claims to close the gap, because
no observation channel built around trade events can.

The fungibility-collapse correspondence does \emph{not} empty
$\ResS$.  $\Pstd$ removes cooperative structure, individuation, and
non-fungibility, but it does not extend the channel's reach;
capabilities exercised without sacrifice events remain unobservable
under either framework.  The residual class is a consequence of the
observation channel, not of the poset structure.  GFM's
contribution beyond standard economics on this front is formal
identification: $\ResS$ is a definable structural object with a
reserved cell ($\delta_{\mathrm{residual}}$) in the gap
decomposition (\citep[Need~Sufficiency]{lasser2026paper8b},
Definition~8) and an explicit governance question (which
capabilities admit ``no sacrifice-observable path in principle?'').
Recognizing the residual class as analytically distinct from
within-channel-but-unwitnessed capabilities (the
$\delta_{\mathrm{dormant}}$ and attestation-limited contributions
of \citep[Need~Sufficiency]{lasser2026paper8b},
Section~5) is an analytical separation that the standard apparatus
collapses.

\begin{remark}[The five features compose]
\label{rem:features_compose}
The five features of this section are not independent: cooperative
capabilities require non-fungibility (cooperative excess only has
content when capabilities are distinct); anti-monopolar reasoning
requires non-fungibility (no dimension to penalize concentration in
on a chain); individuation interacts with cooperative capabilities
through agent-specific bundle compositions; the residual class is
a separate structural object that the other four do not entail.
The composition matters because the structural-prediction class of
Section~\ref{sec:structural_predictions} is generated by
combinations of features rather than by any single feature in
isolation.  Prediction~5 (mixed-polarity bundle-for-bundle traces)
is the canonical example: it requires non-fungibility, cooperative
structure, and individuation simultaneously.
\end{remark}

\subsection{Where GFM and standard economics agree}
\label{sec:fc_agreement}

The correspondence theorem identifies the regime of agreement
precisely.  At the macro scale, with liquid markets, approximately
fungible goods, representative-agent reasoning, and risk-neutral
calibration, standard economics' empirical successes are
inherited by GFM through the correspondence.  At this scale GFM
adds nothing operationally relevant: the cooperative, individuation,
non-fungibility, and anti-monopolar features have content that the
data cannot distinguish from the linear / homogeneous baseline.
The well-known empirical successes of money-maximization
models (consumption smoothing under liquidity constraints,
capital allocation across sectors, intertemporal substitution at
representative-agent scales) are predictions of $\volL$-maximization
on $\Pstd$ as well, and the correspondence shows why.

GFM's added structure becomes load-bearing precisely where the
$\Pstd$ approximation fails:
\begin{itemize}
\item \textbf{Micro scale,} where non-fungibility and cooperative
  structure dominate (consumer bundle decisions, household
  formation, locational choice).
\item \textbf{Non-market sectors,} where the residual-class floor is
  large relative to observed activity (gift economies,
  household production without time-use shadow, volunteer labour
  outside formal channels).
\item \textbf{Multi-substrate domains,} where the anti-monopolar
  property has bite (cross-cultural welfare comparison,
  cross-substrate AI capability landscapes, ecological diversity).
\item \textbf{Heterogeneous-agent populations,} where individuation
  matters and macro-aggregation of preferences is the operational
  question (welfare policy under inequality, cross-cohort
  intertemporal comparison).
\end{itemize}
Each of these regimes generates a divergence from utility-theoretic
predictions that the structural-prediction class of
Section~\ref{sec:structural_predictions} catalogues.  The
fungibility-collapse correspondence is what licenses the claim that
the divergences are GFM-native rather than utility-theoretic
extensions: standard utility theory under fungibility collapse
cannot make these predictions without adding ad hoc structure that
is itself unfalsifiable, while GFM's predictions follow from the
poset-axiomatic apparatus that the correspondence theorem leaves
intact off $\Pstd$.
